\chapter{Especificación protocolos capa de aplicación}
Una de las principales labores a la hora de desarrollar un sistema
distribuido es la definición de las interfaces y del protocolo de
comunicaciones de nivel de aplicación entre las diversas partes que
forman parte del sistema distribuido. El objetivo de este documento no
es otro que mostrar un ejemplo sencillo del diseño y especificación de
una interfaz de comunicaciones y un protocolo de comunicaciones para el
control de forma remota de un coche desde una estación de control.
\subsection{Especificación del protocolo}
Un protocolo queda completamente especificado si definimos su
semántica, sintaxis y temporización, no obstante, para obtener
esta especificación debemos dar una serie de pasos previos. En cualquier caso
el diseño de protocolos es un proceso iterativo que puede dar lugar
a modificaciones a la hora de su implementación.

La especificación de la funcionalidad es el primer paso donde se debe
especificar la funcionalidad esperada del protocolo/interfaz, se deben
describir las partes involucradas y lo que se espera de la interacción
entre las partes. Si como diseñador debe acometer la tarea de realizar una
especificación de la funcionalidad de un protocolo lo ideal es que sea lo menos
ambigua posible, es decir, no esté sujeta a interpretaciones ya que si son
diferentes programadores los que están involucrados en la implementación
del protocolo, una especificación ambigua originará muchos problemas a la hora
de integrar las diferentes partes.

La especificación debería incluir:
%
\begin{itemize}
\item Una descripción genérica explicando la funcionalidad.
\item Los servicios que proveerá y requerirá el protocolo.
\item Especificación de entidades y su relación.
\item Especificación de protocolos entre entidades (Si hay varias entidades usando diversos protocolos).
\item Lista de tipos y formatos de los mensajes intercambiados entre entidades.
\item Temporización describiendo la secuencia de intercambio de mensajes para cada servicio.
\item Definición de interfaces y estructuras de datos en Protobuf.
\end{itemize}

De forma ideal, en protocolos complejos, se debería realizar una descripción formal de cara a eliminar incoherencias y a, en el mejor de los casos, automatizar la generación de las implementaciones y simplificar la verificación.

\begin{figure}[htb]
\centering
\includegraphics[width=0.8\textwidth]{overview.png}
\caption{Protocolo a implementar} \label{fig:overview}
\end{figure}

\subsection{Descripción genérica}
Se trata de diseñar un protocolo de control, es decir, defina
sintaxis, semántica y temporización para el control de un \emph{Unmanned
Ground Vehicle (UGV)} de forma que la comunicación entre la \emph{Estación
Base (EB)} y el vehículo permita monitorizarlo y controlarlo a través
de la red de carreteras nacional. La estación base es un
dispositivo de cómputo con los periféricos adecuados (por ejemplo,
pantalla, teclado, volante/joystick, pedales, etc. ) que permitirá
controlar un UGV a distancia.
\begin{itemize}
\item Puede asumir que su protocolo se
encapsulará sobre una conexión TCP/IP entre cada UGV y la EB.
\item Cada UGV tendrá una dirección IP estática, única y conocida.
\end{itemize}
Para simplificar esta descripción el UGV es un coche de marchas
convencional y asumimos que el lector conoce cómo se conduce un coche
de este tipo de acuerdo a las normas internacionales de circulación.
